\section{候选任务、资源与关联关系}
\label{sec:02_01}

本章考虑一个独立调度窗口内的不可分资源分配问题。决策前已经给定一组候选任务，以及窗口内可供分配的资源。每个候选任务对应一个完整的执行选项，其所需资源必须成组获得：选择该任务，就同时占用其关联的全部资源；不选择该任务，则不产生该任务的资源占用和收益。因此，决策对象是完整任务，而不是任务内部可以任意拆分的资源需求。这里的调度窗口表示统一的资源核算范围，并不包含窗口内部的执行次序。若同一设备可以先后服务多个任务，则其释放与再次占用需要时间信息；仅凭本节的静态关联关系，不能判定这种顺序执行是否可行。

为区分资源身份、任务身份和资源组合，记有限非空资源集合为 $V$，有限非空候选任务索引集合为 $E$。资源标识 $v\in V$ 和任务标识 $e\in E$ 均为无量纲标签，不表示数量或大小。任务 $e$ 所需的非空资源组合记为 $R_e\subseteq V$、$R_e\ne\varnothing$，于是任务与资源之间的基本对应关系为
\begin{equation}
 e\longmapsto R_e\subseteq V,\qquad e\in E.
 \label{s0201:resource_mapping}
\end{equation}
在超图记法 $H=(V,E)$ 中，$E$ 应理解为带任务身份的超边索引族，每条超边的关联资源由 $R_e$ 给出，而不是对资源子集去重后得到的集合。具体而言，两个不同任务 $e,f\in E$ 即使满足 $e\ne f$ 和 $R_e=R_f$，仍然保留各自的任务身份；相同资源组合不意味着相同任务，也不保证二者收益相同。若直接以资源子集替代任务身份并去重，就可能丢失候选选项。沿用本章实验的规模约定，$n=|E|$ 表示候选任务数，$m=|V|$ 表示资源数，二者均为无量纲计数。

一个任务可以关联多个资源，一个资源也可以被多个候选任务共同需要。为准确记录这两种方向的关系，定义关联集合 $\mathcal I$、关联矩阵 $A$ 和资源的关联任务集合 $N(v)$：
\begin{align}
 \mathcal I
 &=\{(v,e)\in V\times E:\ v\in R_e\},
 \label{s0201:incidence_set}\\
 A_{ve}
 &=\begin{cases}
 1,&v\in R_e,\\
 0,&v\notin R_e,
 \end{cases}
 &&\forall v\in V,\ \forall e\in E,
 \label{s0201:incidence_matrix}\\
 N(v)
 &=\{e\in E:\ v\in R_e\},
 &&\forall v\in V.
 \label{s0201:resource_neighborhood}
\end{align}
矩阵 $A\in\{0,1\}^{m\times n}$ 的行对应资源，列对应任务，其元素为无量纲指示量。第 $e$ 列回答任务 $e$ 需要哪些资源，第 $v$ 行回答哪些候选任务会占用资源 $v$。因此，$|R_e|$ 是任务 $e$ 涉及的资源标识数，$|N(v)|$ 是需要资源 $v$ 的候选任务数，均为无量纲计数。二者描述的是关联结构，不能直接解释为执行时长、能耗或实际发生的资源占用量。

任务是否被选中用二元变量 $x_e\in\{0,1\}$ 表示，其中 $x_e=1$ 表示完整选择任务 $e$，$x_e=0$ 表示不选择。本章采用单位需求约定：每个被选任务对其关联的每一种资源恰好占用一份，对未关联资源不产生占用。由此，给定选择向量 $x=(x_e)_{e\in E}\in\{0,1\}^{n}$ 后，资源 $v$ 的占用份数为
\begin{equation}
 \ell_v(x)
 =\sum_{e\in E}A_{ve}x_e
 =\sum_{e\in N(v)}x_e,
 \qquad \forall v\in V.
 \label{s0201:resource_load}
\end{equation}
这里 $\ell_v(x)\in\mathbb Z_{\geq 0}$，按资源 $v$ 自身的可分配单位计数，单位为份。式\eqref{s0201:resource_load} 将同一个 $x_e$ 同时用于任务 $e$ 的全部关联资源，因此不能出现只为该任务分配部分关联资源却仍将任务记为选中的情况。不同资源的份数分别核算，不把不同类型资源的份数相加解释为某个物理总量。若一个任务需要同一资源的多份供给，就必须另行定义整数需求系数及其单位；二元关联量 $A_{ve}$ 本身只回答是否关联，不能同时承担多份需求的含义。

资源 $v$ 在该窗口内的可用容量记为 $b_v\in\mathbb Z_{\geq 0}$，与 $\ell_v(x)$ 使用相同的份数单位。容量是决策前给定的外生参数，用于限制窗口内允许累计分配的份数，而不是资源在未建模的时间顺序中能够重复使用的次数。若 $b_v=0$，任何需要该资源的任务均无法获得完整资源组合；若 $b_v=1$，两个共同需要该资源的任务不能同时入选；容量更大时，则需要检查所有相关入选任务的累计占用。两个任务是否共享资源，可由 $R_e\cap R_f$ 是否为空判断，但共享关系本身不一定意味着互斥，互斥还取决于容量。本书现有纯组合实验全部采用单位容量，因此共享任一资源就构成任务间的两两冲突；这一事实也说明，仅凭任务关联的资源数较多，不能断言模型具有不可约的高阶依赖。

每个候选任务另有一个决策前固定的正权重 $w_e\in\mathbb R_{>0}$，表示完整选择该任务获得的组合收益。本章将其作为无量纲权重使用；若在具体应用中赋予货币或其他收益单位，则各任务权重必须采用统一、可相加的单位。任务 $e$ 的收益贡献为 $w_ex_e$，其含义依赖于完整选择，不能把部分资源分配折算为同一任务的部分收益。固定权重还意味着，本层模型不根据其他任务的选择结果重新计算该任务收益。由此，任务身份、资源组合、可用容量和收益分别承担不同作用：任务身份确定独立决策项，资源组合确定占用去向，容量确定共享资源的承载界限，权重确定选择收益。这些对象共同构成后续容量约束与目标函数的输入，而任务持续时间、资源释放时刻和跨时段复用关系仍需通过新增时间索引及相应状态关系才能表达。
